\documentclass[../DoAn.tex]{subfiles}
\begin{document}

Chương~\ref{chapter:Introduction} đã giới thiệu bài toán gợi ý nhạc Việt theo cảm xúc và theo màu sắc. Chương này khảo sát hiện trạng các hệ thống và nghiên cứu liên quan, rút ra khoảng trống mà đồ án nhắm tới, sau đó phân tích yêu cầu của hệ thống Brightify dưới dạng các tác nhân, các use case chính kèm đặc tả chi tiết, và các yêu cầu chức năng cũng như phi chức năng.

\section{Khảo sát hiện trạng}
\label{section:2.1}
Khảo sát được tiến hành trên ba luồng bổ trợ cho nhau: các hệ thống gợi ý nhạc đang vận hành, các nghiên cứu nhận dạng cảm xúc trong âm nhạc, và các nghiên cứu liên kết màu sắc với âm nhạc. Mỗi luồng làm rõ một mặt của bài toán và cùng nhau chỉ ra khoảng trống mà đồ án nhắm tới.

\subsection{Các hệ thống gợi ý nhạc đang vận hành}
\label{subsection:2.1.1}
Các nền tảng nhạc số phổ biến như Spotify, Apple Music hay các dịch vụ trong nước chủ yếu dựa trên lọc cộng tác (collaborative filtering) và siêu dữ liệu (nghệ sĩ, thể loại, lịch sử nghe). Nhìn rộng hơn, các hệ gợi ý thường được phân thành ba nhóm \cite{aggarwal2016recsys}: \textit{lọc cộng tác} học từ hành vi người dùng, \textit{gợi ý dựa trên nội dung} (content-based) phân tích trực tiếp đặc trưng của vật phẩm, và hệ \textit{lai} (hybrid) kết hợp cả hai \cite{burke2002hybrid} (Hình~\ref{fig:recsys-taxonomy}). Ưu điểm của hướng lọc cộng tác là cá nhân hóa tốt khi có nhiều dữ liệu tương tác: hệ thống học từ hành vi của hàng triệu người dùng để dự đoán sở thích. Tuy nhiên hướng này có ba nhược điểm cố hữu với bài toán của đồ án. Thứ nhất, nó \textit{phụ thuộc dữ liệu người dùng}, vốn nhạy cảm về quyền riêng tư và không phải lúc nào cũng sẵn có. Thứ hai, nó gặp vấn đề \textit{khởi đầu nguội} (cold-start) \cite{schein2002coldstart} với bài hát hoặc người dùng mới chưa có lịch sử tương tác. Thứ ba, và quan trọng nhất, nó \textit{không cho phép truy vấn trực tiếp theo cảm xúc hay theo màu sắc}: người nghe không thể nói ``cho tôi nhạc mang sắc xanh trầm'' mà chỉ có thể chờ hệ thống suy ra sở thích từ hành vi quá khứ. Đây chính là khoảng trống về trải nghiệm mà đồ án muốn lấp; Brightify chọn hướng \textit{dựa trên nội dung} để vừa cho phép truy vấn theo cảm xúc/màu, vừa tránh phụ thuộc dữ liệu hành vi và bài toán khởi đầu nguội. Bảng~\ref{table:vnplatforms} đối chiếu các nền tảng nhạc số phổ biến với Brightify theo khả năng truy vấn cảm xúc/màu và mức phụ thuộc dữ liệu người dùng.

\begin{figure}[tbp]
\centering
\begin{tikzpicture}[font=\small,
   box/.style={draw, rounded corners, align=center, minimum width=2.7cm, minimum height=0.85cm, fill=black!3},
   root/.style={draw, rounded corners, align=center, minimum width=2.6cm, minimum height=0.85cm, fill=blue!8},
   pick/.style={draw, thick, rounded corners, align=center, minimum width=2.7cm, minimum height=0.85cm, fill=green!10},
   arr/.style={-{Stealth[length=2mm]}, thick}]
  \node[root] (r) {Hệ gợi ý};
  \node[box] (cf) at (-3.5,-1.6) {Lọc cộng tác\\\footnotesize(collaborative)};
  \node[pick] (cb) at (0,-1.6) {Dựa trên nội dung\\\footnotesize(content-based)};
  \node[box] (hy) at (3.5,-1.6) {Lai\\\footnotesize(hybrid)};
  \draw[arr] (r) -- (cf); \draw[arr] (r) -- (cb); \draw[arr] (r) -- (hy);
  \node[font=\scriptsize, red!60!black, align=center] at (-3.5,-2.7) {dễ gặp khởi đầu nguội};
  \node[font=\scriptsize, green!45!black, align=center] at (0,-2.7) {miễn nhiễm khởi đầu nguội\\\textbf{(Brightify)}};
  \node[font=\scriptsize, align=center] at (3.5,-2.7) {kết hợp cả hai};
\end{tikzpicture}
\caption{Phân loại hệ gợi ý: lọc cộng tác phụ thuộc hành vi người dùng nên dễ gặp khởi đầu nguội; Brightify đi theo hướng dựa trên nội dung}
\label{fig:recsys-taxonomy}
\end{figure}

\begin{longtable}{|p{3.0cm}|p{4.0cm}|p{1.9cm}|p{1.5cm}|p{2.0cm}|}
\caption{Khảo sát khả năng truy vấn theo cảm xúc/màu của các nền tảng nhạc số (đối chiếu định tính)}\label{table:vnplatforms}\\
\hline
\textbf{Nền tảng} & \textbf{Cơ chế gợi ý chính} & \textbf{Truy vấn cảm xúc} & \textbf{Truy vấn màu} & \textbf{Phụ thuộc dữ liệu người dùng} \\
\hline
\endfirsthead
\hline
\textbf{Nền tảng} & \textbf{Cơ chế gợi ý chính} & \textbf{Truy vấn cảm xúc} & \textbf{Truy vấn màu} & \textbf{Phụ thuộc dữ liệu người dùng} \\
\hline
\endhead
Spotify, Apple Music & Lọc cộng tác + biên tập + đặc trưng âm thanh & Danh sách phát tâm trạng biên tập & Không & Cao \\
\hline
YouTube Music & Lọc cộng tác + lịch sử xem & Hạn chế & Không & Cao \\
\hline
Zing MP3, NhacCuaTui & Lọc cộng tác + bảng xếp hạng + biên tập & Hạn chế & Không & Cao \\
\hline
\textbf{Brightify} & \textbf{Dựa trên nội dung trên trục V-A} & \textbf{Có (theo cảm xúc)} & \textbf{Có} & \textbf{Không} \\
\hline
\end{longtable}

Bài toán khởi đầu nguội thực ra có ba dạng \cite{schein2002coldstart}: \textit{người dùng mới} (chưa có lịch sử để suy sở thích), \textit{vật phẩm mới} (bài hát mới chưa ai nghe nên lọc cộng tác không có tín hiệu tương tác), và \textit{hệ thống mới} (toàn bộ nền tảng chưa tích lũy dữ liệu hành vi). Lọc cộng tác bất lực với cả ba dạng vì nó hoàn toàn dựa trên ma trận tương tác người dùng--vật phẩm. Ngược lại, gợi ý dựa trên nội dung miễn nhiễm với khởi đầu nguội \textit{vật phẩm} và \textit{hệ thống}, vì nó chấm điểm dựa trên đặc trưng nội tại của bài hát (âm thanh, lời, cảm xúc) chứ không cần bất kỳ lịch sử nghe nào --- một bài vừa được thêm vào kho có thể được gợi ý ngay sau khi đặc trưng của nó được trích. Đây là lý do nền tảng để đồ án chọn hướng nội dung: nó phù hợp với bối cảnh không có dữ liệu người dùng và một kho nhạc tự thu thập có thể liên tục mở rộng.

Cần khách quan nhìn nhận rằng hướng dựa trên nội dung cũng có những đánh đổi riêng. Vì chỉ dựa vào đặc trưng nội tại của bài hát, nó dễ \textit{quá chuyên biệt} (over-specialization): kết quả có xu hướng giống bài gốc và ít bất ngờ thú vị (serendipity) hơn so với lọc cộng tác, vốn khai thác được các liên kết kiểu ``người thích A cũng thích B'' bất kể A và B có giống nhau về nội dung hay không. Ngoài ra, chất lượng gợi ý phụ thuộc trực tiếp vào chất lượng đặc trưng trích được: nếu biểu diễn âm thanh hay tín hiệu cảm xúc nghèo nàn thì kết quả cũng kém theo. Đồ án chấp nhận các đánh đổi này vì chúng phù hợp với mục tiêu (cho phép truy vấn theo cảm xúc và màu, không cần dữ liệu người dùng), và bù lại bằng cách đầu tư mạnh vào chất lượng đặc trưng --- mô hình tự giám sát mạnh cho âm thanh, tổ hợp nhiều tín hiệu cho cảm xúc --- đúng phần mà hướng nội dung có thể kiểm soát được. Sự đánh đổi giữa các hướng tiếp cận được tóm tắt ở Bảng~\ref{table:khaosat} phía dưới.

\subsection{Nghiên cứu nhận dạng cảm xúc trong âm nhạc}
\label{subsection:2.1.2}
Bài toán nhận dạng cảm xúc trong âm nhạc (Music Emotion Recognition, MER) \cite{yang2012mer} đã được nghiên cứu nhiều trên các tập dữ liệu tiếng Anh được con người gán nhãn như DEAM \cite{aljanaki2017deam} và PMEmo \cite{zhang2018pmemo}, thường biểu diễn cảm xúc trên không gian Valence--Arousal \cite{russell1980circumplex}. Các nghiên cứu này cho thấy Arousal (mức năng lượng) dễ dự đoán từ đặc trưng âm thanh hơn Valence (mức tích cực), vì Valence gắn nhiều với nội dung ngữ nghĩa của lời hơn là với âm sắc. Khảo sát xuyên văn hóa gần đây cho thấy cấu trúc Valence--Arousal khá ổn định giữa các nền văn hóa, dù các nhãn cảm xúc cụ thể vẫn phân kỳ theo ngôn ngữ \cite{lee2025globalmood} --- điều này vừa ủng hộ việc dùng V-A làm trục chung, vừa cho thấy cần một xử lý riêng cho tiếng Việt. Về phương pháp, các nghiên cứu MER ban đầu dùng đặc trưng âm thanh thủ công (MFCC, chroma, tempo, đặc trưng phổ) kết hợp hồi quy hoặc phân loại; gần đây chuyển sang đặc trưng học sâu, đặc biệt là các mô hình học tự giám sát trên lượng lớn nhạc không nhãn (xem Chương~\ref{chapter:Methodology}). Một phát hiện nhất quán xuyên các nghiên cứu là \textit{Valence khó dự đoán từ âm thanh hơn Arousal}: Arousal gắn chặt với các thuộc tính âm thanh đo được (nhịp độ, độ to, năng lượng phổ) nên dễ hồi quy, trong khi Valence phụ thuộc nhiều vào ngữ nghĩa và bối cảnh văn hóa của lời, vốn không hiện diện đầy đủ trong tín hiệu âm thanh. Quan sát này định hướng trực tiếp thiết kế của đồ án: tách hai trục và suy Valence chủ yếu từ lời, Arousal chủ yếu từ âm thanh (mục~\ref{subsection:4.2.1}).

Về mặt phương pháp, lĩnh vực MER đã trải qua một sự dịch chuyển đáng chú ý. Thế hệ đầu dựa vào \textit{đặc trưng âm thanh thủ công} (MFCC mô tả âm sắc, chroma mô tả hòa âm, cùng các thống kê về nhịp độ, độ to và năng lượng phổ) đưa vào các mô hình hồi quy hoặc phân loại cổ điển; cách này minh bạch nhưng bị giới hạn bởi chất lượng của đặc trưng được thiết kế tay. Thế hệ tiếp theo dùng mạng học sâu huấn luyện đầu cuối, đạt độ chính xác cao hơn nhưng đòi hỏi lượng lớn dữ liệu gán nhãn --- điều khó đáp ứng với MER, nơi nhãn cảm xúc do con người chấm rất tốn kém. Gần đây nhất, hướng \textit{học tự giám sát} tiền huấn luyện trên lượng lớn nhạc không nhãn rồi gắn một đầu dò nhẹ phía trên đã trở thành cách tiếp cận chủ đạo, vì nó kế thừa được sức mạnh biểu diễn của học sâu mà không cần nhiều nhãn. Đồ án đi theo đúng hướng hiện đại này: dùng một mô hình tự giám sát đóng băng làm bộ trích đặc trưng và chỉ huấn luyện một đầu dò tuyến tính (chi tiết ở Chương~\ref{chapter:Methodology}), một lựa chọn đặc biệt phù hợp khi dữ liệu nhãn khan hiếm như với nhạc Việt.

Đáng chú ý, trong phạm vi khảo sát của đồ án, chưa tìm thấy tập dữ liệu cảm xúc âm nhạc \textit{tiếng Việt} được con người gán nhãn và công bố công khai; các tài nguyên cảm xúc tiếng Việt hiện có chủ yếu là văn bản mạng xã hội, không phải nhạc. Đây là khoảng trống về dữ liệu buộc đồ án phải tìm cách gán cảm xúc mà không dựa vào nhãn người bản địa. Bảng~\ref{table:merdata} tóm tắt các tập dữ liệu cảm xúc nền mà đồ án dùng để hiệu chỉnh và kiểm chứng, đồng thời cho thấy rõ khoảng trống dữ liệu tiếng Việt.

\begin{longtable}{|p{2.6cm}|p{2.6cm}|p{2.0cm}|p{2.2cm}|p{3.0cm}|}
\caption{Các tập dữ liệu cảm xúc nền (gán nhãn người) và vai trò trong đồ án}\label{table:merdata}\\
\hline
\textbf{Tập dữ liệu} & \textbf{Loại dữ liệu} & \textbf{Cỡ} & \textbf{Nhãn} & \textbf{Vai trò trong đồ án} \\
\hline
\endfirsthead
\hline
\textbf{Tập dữ liệu} & \textbf{Loại dữ liệu} & \textbf{Cỡ} & \textbf{Nhãn} & \textbf{Vai trò trong đồ án} \\
\hline
\endhead
DEAM \cite{aljanaki2017deam} & Nhạc (Tây) & 1.802 bài & V-A liên tục & Huấn luyện đầu dò Arousal \\
\hline
PMEmo \cite{zhang2018pmemo} & Nhạc, đoạn chorus (Tây) & 794 bài & V-A động & Kiểm chứng chuyển miền \\
\hline
EmoBank \cite{buechel2017emobank} & Văn bản (Anh) & $\approx$10.000 câu & V-A-D & Huấn luyện đầu dò Valence liên ngữ \\
\hline
Nhạc Việt gán nhãn & Nhạc (Việt) & \textit{không tồn tại} & --- & Khoảng trống $\Rightarrow$ EWE không nhãn \\
\hline
\end{longtable}

\subsection{Liên kết màu sắc và âm nhạc}
\label{subsection:2.1.3}
Hướng liên kết màu--âm nhạc không phải ý tưởng hoàn toàn mới. Quan hệ màu--nhạc ở con người đã được chứng minh là \textit{trung gian hóa bởi cảm xúc} với mức tương quan rất cao \cite{palmer2013music, whiteford2018color}, dựa trên các chuẩn liên kết màu--cảm xúc xuyên văn hóa \cite{jonauskaite2020universal}; rộng hơn, đây là một trường hợp của \textit{tương ứng đa giác quan} (crossmodal correspondence) giữa thị giác và thính giác \cite{spence2011crossmodal}; đã có tập dữ liệu gắn màu, cảm xúc và nhạc trên cùng mô hình V-A phục vụ truy hồi theo cảm xúc \cite{pesek2017moodo}; và đã có hệ thống gợi ý nhạc từ màu \cite{kittimathaveenan2020color}. Cơ chế đằng sau liên kết màu--nhạc đáng được nhấn mạnh: nghiên cứu của Palmer và cộng sự \cite{palmer2013music} cho thấy khi người nghe được yêu cầu chọn màu ``hợp'' với một đoạn nhạc, lựa chọn của họ tương quan rất cao với \textit{cảm xúc} mà cả màu lẫn nhạc gợi lên --- nghĩa là cảm xúc là biến trung gian, không phải một liên kết màu--nhạc trực tiếp. Whiteford và cộng sự \cite{whiteford2018color} mở rộng quan sát này thành quan hệ định lượng giữa thuộc tính màu (độ sáng, độ bão hòa) và thuộc tính nhạc (nhịp độ, năng lượng). Đây chính là cơ sở để đồ án không ghép màu với nhạc trực tiếp mà ghép \textit{qua} không gian cảm xúc Valence--Arousal. Tuy nhiên, các công trình ứng dụng trước đây gắn với dữ liệu phương Tây và thường ánh xạ màu sang thuộc tính nhạc bằng \textit{quy tắc thủ công}, chưa dùng một không gian cảm xúc \textit{học được} dùng chung cho cả hai chức năng, và chưa nhắm tới nhạc Việt. Bảng~\ref{table:colorprior} đối chiếu các công trình màu--nhạc tiêu biểu với hệ thống đề xuất theo bốn tiêu chí kỹ thuật.

\begin{longtable}{|p{3.0cm}|p{2.4cm}|p{3.0cm}|p{2.6cm}|}
\caption{So sánh các công trình màu$\rightarrow$nhạc tiêu biểu với hệ thống đề xuất}\label{table:colorprior}\\
\hline
\textbf{Công trình} & \textbf{Không gian màu} & \textbf{Cách ánh xạ màu$\rightarrow$nhạc} & \textbf{V-A dùng chung} \\
\hline
\endfirsthead
\hline
\textbf{Công trình} & \textbf{Không gian màu} & \textbf{Cách ánh xạ màu$\rightarrow$nhạc} & \textbf{V-A dùng chung} \\
\hline
\endhead
MOODO \cite{pesek2017moodo} & RGB/HSV & Gán nhãn người trên tập riêng & Một phần \\
\hline
Kittimathaveenan \cite{kittimathaveenan2020color} & RGB & Quy tắc thủ công & Không \\
\hline
\textbf{Brightify (đề xuất)} & \textbf{Oklab} & \textbf{Hồi quy về chuẩn ICEAS + Whiteford} & \textbf{Có} \\
\hline
\end{longtable}

Bảng~\ref{table:khaosat} tổng hợp ở mức cao hơn, so sánh ba hướng tiếp cận lớn với hệ thống đề xuất. Khoảng trống mà đồ án nhắm tới không phải là ``ý tưởng gợi ý theo màu'' (vốn đã có tiền lệ), mà là một hệ thống gợi ý nhạc \textit{Việt} dựa trên cảm xúc, không phụ thuộc dữ liệu người dùng, hợp nhất truy vấn theo bài hát và theo màu trên một không gian cảm xúc \textit{học được} dùng chung, và có cơ sở khoa học rõ ràng cho từng thành phần.

\begin{longtable}{|p{3.2cm}|p{2.6cm}|p{2.4cm}|p{3.2cm}|}
\caption{So sánh các hướng tiếp cận hiện có và hệ thống đề xuất}\label{table:khaosat}\\
\hline
\textbf{Tiêu chí} & \textbf{Lọc cộng tác (Spotify, v.v.)} & \textbf{Nghiên cứu MER} & \textbf{Brightify (đề xuất)} \\
\hline
\endfirsthead
\hline
\textbf{Tiêu chí} & \textbf{Lọc cộng tác (Spotify, v.v.)} & \textbf{Nghiên cứu MER} & \textbf{Brightify (đề xuất)} \\
\hline
\endhead
Cần dữ liệu người dùng & Có & Không & Không \\
\hline
Gợi ý theo cảm xúc & Gián tiếp & Có (phân tích) & Có \\
\hline
Gợi ý theo màu & Không & Không & Có \\
\hline
Hỗ trợ nhạc Việt & Một phần & Hạn chế & Trọng tâm \\
\hline
\end{longtable}

\subsection{Bối cảnh và thách thức dữ liệu nhạc Việt}
\label{subsection:2.1.4}
Ba luồng khảo sát trên hội tụ về một thách thức chung khi áp dụng cho nhạc Việt: \textit{sự khan hiếm dữ liệu cảm xúc bản địa}. Thị trường nhạc số Việt Nam sôi động với nhiều nền tảng (các dịch vụ trong nước và quốc tế), nhưng dữ liệu cảm xúc gắn với nhạc Việt hầu như không được công bố công khai dưới dạng có thể huấn luyện. Các tài nguyên cảm xúc tiếng Việt sẵn có (như tập cảm xúc văn bản mạng xã hội) thuộc miền văn bản, không phải miền âm nhạc, và cũng không gán theo trục Valence--Arousal liên tục. Hệ quả là không thể huấn luyện trực tiếp một mô hình cảm xúc nhạc Việt theo cách giám sát thông thường.

Thách thức này định hình ba quyết định thiết kế cốt lõi của đồ án. Thứ nhất, vì không có nhãn người bản địa, hệ thống phải gán cảm xúc bằng cách \textit{chuyển giao} từ các nguồn có nhãn (đầu dò tuyến tính trên đặc trưng đóng băng, tổ hợp nhiều tín hiệu) thay vì học giám sát end-to-end. Thứ hai, vì không thu thập dữ liệu hành vi người dùng (do ràng buộc riêng tư và phạm vi), hệ thống không thể dựa vào lọc cộng tác mà phải dựa hoàn toàn vào nội dung. Thứ ba, vì dữ liệu phải tự thu thập, quy trình thu thập (mô tả ở mục~\ref{section:data}) phải tự động và lặp lại được. Những ràng buộc này không phải hạn chế tình thế mà được biến thành nguyên tắc thiết kế nhất quán xuyên suốt đồ án.

\subsection{Đánh giá hệ gợi ý khi không có người dùng}
\label{subsection:2.1.5}
Vì đồ án chủ trương không thu thập dữ liệu người dùng, một câu hỏi phương pháp luận đặt ra ngay từ khâu khảo sát là: \textit{làm thế nào đánh giá một hệ gợi ý khi không có người dùng để đo lường?} Tài liệu về đánh giá hệ gợi ý cung cấp hai nhóm độ đo ngoại tuyến. Nhóm thứ nhất đo độ chính xác xếp hạng so với một tập liên quan (ground truth), tiêu biểu là độ lợi tích lũy chiết khấu chuẩn hóa NDCG \cite{jarvelin2002ndcg}, vốn thưởng cho việc đặt các vật phẩm liên quan lên đầu danh sách. Nhóm thứ hai đo các tính chất \textit{nội tại} của danh sách kết quả mà không cần nhãn liên quan tuyệt đối, chẳng hạn mức danh sách phản ánh đúng phân bố sở thích (calibrated recommendations) \cite{steck2018calibrated}, hay độ nhất quán và độ đa dạng nội bộ của kết quả.

Tuy nhiên, đánh giá ngoại tuyến có hai cạm bẫy đã được ghi nhận rõ. Thứ nhất là \textit{khoảng cách offline--online}: theo Dacrema và cộng sự \cite{dacrema2019progress}, tương quan giữa các độ đo ngoại tuyến và mức hài lòng thực tế của người dùng chỉ vào khoảng $r \approx 0.28$, nên một cải thiện ngoại tuyến không tự động kéo theo cải thiện trải nghiệm. Thứ hai là vấn đề \textit{phương pháp nền yếu}: nhiều công bố báo cáo tiến bộ chỉ vì so với phương pháp nền được cấu hình kém, và Dacrema và cộng sự cho thấy không ít phương pháp ``tân tiến'' không vượt nổi các phương pháp nền đơn giản nếu các phương pháp nền đó được tinh chỉnh tử tế. Hai bài học này định hình trực tiếp quy trình đánh giá của đồ án (Chương~\ref{chapter:Evaluation}): luôn đối chiếu với phương pháp nền mạnh (kể cả phương pháp nền theo độ phổ biến), bổ sung các kiểm chứng \textit{độc lập với nhãn} để giảm phụ thuộc vào một độ đo tự tham chiếu, và nêu thẳng rằng mọi kết quả là độ đo ngoại tuyến chứ chưa phải bằng chứng từ người nghe thật. Như vậy, ràng buộc ``không có người dùng'' không loại bỏ khả năng đánh giá, mà chuyển trọng tâm sang một quy trình ngoại tuyến nhiều tầng, có kiểm định thống kê và tam giác hóa bằng nhiều nguồn độc lập.

\section{Tác nhân và tổng quan chức năng}
\label{section:2.2}

\subsection{Các tác nhân của hệ thống}
\label{subsection:2.2.1}
Khác với một ứng dụng nghe nhạc thương mại thông thường, Brightify gồm hai pha rõ rệt: pha \textit{phục vụ} (serving) tương tác với người nghe theo thời gian thực, và pha \textit{ngoại tuyến} (offline) chuẩn bị toàn bộ đặc trưng cho kho nhạc. Vì vậy mô hình tác nhân của hệ thống gồm bốn tác nhân, trong đó có cả tác nhân chính (con người) và tác nhân phụ (hệ thống/dữ liệu). Thứ nhất, (i) \textbf{người nghe / người dùng} là tác nhân chính của pha phục vụ: họ dùng giao diện để chọn màu hoặc chọn bài hát, xem và phát danh sách bài hát được gợi ý, và xem thông tin cảm xúc của bài hát, và không cần đăng nhập. Thứ hai, (ii) \textbf{quản trị dữ liệu / người vận hành hệ thống} cũng là tác nhân chính, nhưng làm việc với pha ngoại tuyến: họ chịu trách nhiệm cập nhật kho nhạc (thêm bài mới, lọc trùng/cover) và kích hoạt việc chạy lại quy trình trích đặc trưng ngoại tuyến. Bên cạnh hai tác nhân người, hệ thống có hai tác nhân phụ: (iii) \textbf{kho nhạc / dataset} là tập dữ liệu nhạc Việt gồm tệp âm thanh, siêu dữ liệu và lời bài hát, cung cấp dữ liệu đầu vào cho các use case phục vụ và là đối tượng được cập nhật bởi người vận hành; và (iv) \textbf{hệ thống xử lý đặc trưng offline} tự động trích embedding âm thanh (MuQ), tính tín hiệu cảm xúc từ lời, nhịp độ, độ to và sinh các ma trận biểu diễn, được người vận hành kích hoạt và ghi kết quả trở lại kho đặc trưng để pha phục vụ nạp lên.

Việc tách hai tác nhân phụ (kho nhạc và hệ thống xử lý đặc trưng offline) là một lựa chọn thiết kế có chủ đích: nó phản ánh đúng kiến trúc hai pha của hệ thống (Chương~\ref{chapter:Design}), trong đó toàn bộ tính toán nặng được dồn về pha ngoại tuyến để pha phục vụ chỉ còn các phép tra cứu và nhân ma trận. Bảng~\ref{table:actors} tóm tắt trách nhiệm chính của từng tác nhân.

\begin{longtable}{|p{4.0cm}|p{3.0cm}|p{6.0cm}|}
\caption{Trách nhiệm chính của các tác nhân}\label{table:actors}\\
\hline
\textbf{Tác nhân} & \textbf{Loại} & \textbf{Trách nhiệm chính} \\
\hline
\endfirsthead
\hline
\textbf{Tác nhân} & \textbf{Loại} & \textbf{Trách nhiệm chính} \\
\hline
\endhead
Người nghe & Chính (người) & Chọn màu/bài, xem và phát danh sách phát, xem cảm xúc bài hát \\
\hline
Người vận hành & Chính (người) & Cập nhật kho nhạc, kích hoạt pipeline ngoại tuyến \\
\hline
Kho nhạc / dataset & Phụ (dữ liệu) & Cung cấp âm thanh, siêu dữ liệu, lời; là đối tượng cập nhật \\
\hline
Hệ thống xử lý offline & Phụ (hệ thống) & Trích đặc trưng, tính tọa độ cảm xúc, sinh ma trận biểu diễn \\
\hline
\end{longtable}

\subsection{Biểu đồ use case tổng quát}
\label{subsection:2.2.2}
Hình~\ref{fig:usecase} trình bày biểu đồ use case tổng quát của hệ thống. Người nghe thao tác với bốn use case của pha phục vụ; người vận hành thao tác với hai use case của pha ngoại tuyến; hai tác nhân phụ (kho nhạc và hệ thống xử lý đặc trưng offline) tham gia cung cấp/nhận dữ liệu cho các use case tương ứng.

\begin{figure}[tbp]
\centering
\tikzset{
  human/.pic={
    \draw[thick] (0,0.55) circle (0.16);
    \draw[thick] (0,0.39) -- (0,-0.12);
    \draw[thick] (-0.22,0.18) -- (0.22,0.18);
    \draw[thick] (0,-0.12) -- (-0.17,-0.48);
    \draw[thick] (0,-0.12) -- (0.17,-0.48);
  }
}
\resizebox{0.92\textwidth}{!}{%
\begin{tikzpicture}[
   uc/.style={draw, ellipse, align=center, minimum width=4.0cm, minimum height=1.3cm},
   sys/.style={draw, rounded corners, align=center, fill=black!4, minimum width=2.9cm, minimum height=1.0cm},
   every node/.style={font=\small}]

  % System boundary
  \draw[thick, rounded corners] (3.4,-0.9) rectangle (10.6,11.7);
  \node[anchor=north] at (7.0,11.5) {\textbf{Hệ thống Brightify}};

  % Use cases — serving group
  \node[uc] (uc1) at (7.0,10.0) {UC01. Chọn bài hát\\để gợi ý bài tương tự};
  \node[uc] (uc2) at (7.0,8.1) {UC02. Chọn màu\\để gợi ý danh sách phát};
  \node[uc] (uc3) at (7.0,6.2) {UC03. Xem danh sách\\bài gợi ý};
  \node[uc] (uc4) at (7.0,4.3) {UC04. Xem thông tin\\cảm xúc của bài hát};
  % Use cases — offline group
  \node[uc] (uc5) at (7.0,2.0) {UC05. Cập nhật\\kho nhạc};
  \node[uc] (uc6) at (7.0,0.1) {UC06. Chạy lại pipeline\\trích đặc trưng};

  % Actors
  \pic at (0.3,6.2) {human}; \node[align=center] at (0.3,5.1) {\textbf{Người nghe}\\(người dùng)};
  \pic at (13.3,1.0) {human}; \node[align=center] at (13.3,-0.1) {\textbf{Người vận hành}\\/ quản trị dữ liệu};
  \node[sys] (store) at (13.3,5.2) {Kho nhạc /\\Dataset};
  \node[sys] (offline) at (13.3,3.3) {Hệ thống xử lý\\đặc trưng offline};

  % Associations — listener fans out to UC01--UC04
  \draw (1.1,5.9) -- (uc1.west);
  \draw (1.1,5.9) -- (uc2.west);
  \draw (1.1,5.9) -- (uc3.west);
  \draw (1.1,5.9) -- (uc4.west);
  % Associations — operator to UC05--UC06
  \draw (12.5,1.0) -- (uc5.east);
  \draw (12.5,1.0) -- (uc6.east);
  % Secondary actors (dashed)
  \draw[dashed] (store.west) -- (uc2.east);
  \draw[dashed] (store.west) -- (uc5.east);
  \draw[dashed] (offline.west) -- (uc6.east);
\end{tikzpicture}}
\caption{Biểu đồ use case tổng quát của hệ thống Brightify}
\label{fig:usecase}
\end{figure}

Bảng~\ref{table:uclist} liệt kê các use case và tác nhân liên quan. Sáu use case chia thành hai nhóm: nhóm phục vụ người nghe (UC01--UC04) và nhóm vận hành dữ liệu (UC05--UC06). Một số use case phụ trợ (tìm kiếm bài hát, phát nhạc và quản lý hàng đợi) được đặc tả trong Phụ lục~\ref{appendix:usecase}.

\begin{longtable}{|p{1.2cm}|p{6.0cm}|p{6.0cm}|}
\caption{Danh sách use case của hệ thống}\label{table:uclist}\\
\hline
\textbf{Mã} & \textbf{Use case} & \textbf{Tác nhân} \\
\hline
\endfirsthead
\hline
\textbf{Mã} & \textbf{Use case} & \textbf{Tác nhân} \\
\hline
\endhead
UC01 & Chọn bài hát để gợi ý bài tương tự & Người nghe (chính); kho nhạc (phụ) \\
\hline
UC02 & Chọn màu để gợi ý danh sách phát & Người nghe (chính); kho nhạc (phụ) \\
\hline
UC03 & Xem danh sách bài gợi ý & Người nghe \\
\hline
UC04 & Xem thông tin cảm xúc của bài hát & Người nghe \\
\hline
UC05 & Cập nhật kho nhạc & Người vận hành (chính); kho nhạc (phụ) \\
\hline
UC06 & Chạy lại pipeline trích đặc trưng & Người vận hành (chính); hệ thống xử lý đặc trưng offline (phụ) \\
\hline
\end{longtable}

\section{Yêu cầu chức năng}
\label{section:2.fr}
Từ các tác nhân và use case ở trên, đồ án rút ra các yêu cầu chức năng (functional requirements) cụ thể mà hệ thống phải đáp ứng, được liệt kê trong Bảng~\ref{table:fr} kèm mức độ ưu tiên và use case liên quan. Các yêu cầu được phân thành ba nhóm: nhóm gợi ý cốt lõi (FR01--FR03) là trọng tâm khoa học của đồ án; nhóm tinh chỉnh chất lượng kết quả (FR04--FR06) bảo đảm danh sách trả về thực sự dùng được; và nhóm hỗ trợ trải nghiệm và vận hành (FR07--FR10).

\begin{longtable}{|p{1.0cm}|p{7.6cm}|p{2.2cm}|p{1.8cm}|}
\caption{Danh sách yêu cầu chức năng của hệ thống}\label{table:fr}\\
\hline
\textbf{Mã} & \textbf{Yêu cầu chức năng} & \textbf{Use case} & \textbf{Ưu tiên} \\
\hline
\endfirsthead
\hline
\textbf{Mã} & \textbf{Yêu cầu chức năng} & \textbf{Use case} & \textbf{Ưu tiên} \\
\hline
\endhead
FR01 & Gợi ý các bài hát tương tự một bài cho trước (âm thanh + cảm xúc + lời) & UC01 & Cao \\
\hline
FR02 & Gợi ý danh sách phát từ một màu sắc người dùng chọn & UC02 & Cao \\
\hline
FR03 & Tạo ``hành trình cảm xúc'' từ hai màu theo nguyên lý đồng cảm & UC02 & Cao \\
\hline
FR04 & Loại bỏ bản cover/trùng khỏi kết quả gợi ý tương tự & UC01 & Trung bình \\
\hline
FR05 & Đa dạng hóa kết quả theo nghệ sĩ (tránh lặp nghệ sĩ) & UC01, UC02 & Trung bình \\
\hline
FR06 & Hỗ trợ chế độ ``đài'' vô tận, không lặp bài đã phát & UC03 & Trung bình \\
\hline
FR07 & Hiển thị thông tin cảm xúc của bài hát (V-A, nhãn tâm trạng, lời) & UC04 & Cao \\
\hline
FR08 & Duyệt, lọc và sắp xếp kho nhạc & UC03 & Trung bình \\
\hline
FR09 & Tìm kiếm bài hát theo tên, nghệ sĩ, lời hoặc ``cảm giác'' & --- & Thấp \\
\hline
FR10 & Cập nhật kho nhạc và chạy lại pipeline trích đặc trưng & UC05, UC06 & Trung bình \\
\hline
\end{longtable}

Trong các yêu cầu trên, FR01 và FR02 là hai chức năng cốt lõi đã nêu ở Chương~\ref{chapter:Introduction}; FR03 mở rộng FR02 cho trường hợp hai màu. FR04--FR06 không phải ``tính năng phụ'' mà là điều kiện để kết quả gợi ý có chất lượng: nếu không loại cover thì danh sách tương tự sẽ trả về chính bài nguồn dưới dạng bản phối khác, còn nếu không đa dạng nghệ sĩ thì danh sách phát dễ bị một nghệ sĩ chiếm trọn. FR07 hiện thực hóa tính ``giải thích được'' của hệ thống. FR08--FR10 bảo đảm hệ thống vận hành được như một sản phẩm hoàn chỉnh.

\section{Đặc tả chi tiết use case}
\label{section:2.3}
Phần này đặc tả các use case quan trọng theo bốn yếu tố: tiền điều kiện, luồng sự kiện chính, luồng thay thế/ngoại lệ và hậu điều kiện. Đặc tả của sáu use case nghiệp vụ chính lần lượt được trình bày trong Bảng~\ref{table:uc01} (UC01), Bảng~\ref{table:uc02} (UC02), Bảng~\ref{table:uc03} (UC03), Bảng~\ref{table:uc04} (UC04), Bảng~\ref{table:uc05} (UC05) và Bảng~\ref{table:uc06} (UC06). Các con số kỹ thuật (cấu trúc dữ liệu trả về, mã lỗi) được trình bày trong Chương~\ref{chapter:Implementation}, còn thời gian phản hồi đo được trình bày trong Chương~\ref{chapter:Evaluation}.

\begin{longtable}{|p{3.0cm}|p{11.0cm}|}
\caption{Đặc tả use case UC01 --- Gợi ý bài tương tự}\label{table:uc01}\\
\hline
\textbf{Mã / Tên} & UC01 --- Chọn bài hát để gợi ý bài tương tự \\
\hline
\endfirsthead
\hline
\textbf{Mã / Tên} & UC01 --- Chọn bài hát để gợi ý bài tương tự \\
\hline
\endhead
\textbf{Tác nhân} & Người nghe (chính); kho nhạc và đặc trưng đã tính sẵn (phụ) \\ \hline
\textbf{Mô tả} & Từ một bài hát đang nghe hoặc được chọn, hệ thống trả về danh sách các bài ``nghe giống'' về âm thanh và gần về cảm xúc. \\ \hline
\textbf{Tiền điều kiện} & Kho nhạc và các đặc trưng (embedding âm thanh MuQ, tọa độ cảm xúc, embedding lời) đã được nạp; bài hát nguồn tồn tại trong kho. \\ \hline
\textbf{Luồng sự kiện chính} & 1. Người nghe chọn một bài hát (từ thư viện, kết quả tìm kiếm hoặc bài đang phát). \newline 2. Hệ thống tính điểm tương đồng giữa bài nguồn và mọi bài khác bằng tổ hợp ba tín hiệu (âm thanh, cảm xúc, lời). \newline 3. Hệ thống loại các bản trùng/cover của bài nguồn và đa dạng hóa theo nghệ sĩ. \newline 4. Hệ thống trả về danh sách $top\text{-}k$ bài tương tự kèm điểm tương đồng. \newline 5. Người nghe xem danh sách (chuyển sang UC03) và có thể phát. \\ \hline
\textbf{Luồng thay thế / ngoại lệ} & 2a. Bài nguồn không có đặc trưng âm thanh: hệ thống lùi về dùng tín hiệu cảm xúc/lời còn lại. \newline 3a. Người nghe yêu cầu thêm kết quả (chế độ ``đài'' vô tận): hệ thống nối tiếp danh sách và loại các bài đã phát qua tham số loại trừ. \newline 1a. Không xác định được bài nguồn: hệ thống báo lỗi ``không tìm thấy bài hát''. \\ \hline
\textbf{Hậu điều kiện} & Một danh sách bài tương tự (không trùng bài nguồn, đã đa dạng nghệ sĩ) được hiển thị; trạng thái hàng đợi phát được cập nhật nếu người nghe chọn phát. \\ \hline
\end{longtable}

\begin{longtable}{|p{3.0cm}|p{11.0cm}|}
\caption{Đặc tả use case UC02 --- Gợi ý theo màu}\label{table:uc02}\\
\hline
\textbf{Mã / Tên} & UC02 --- Chọn màu để gợi ý danh sách phát \\
\hline
\endfirsthead
\hline
\textbf{Mã / Tên} & UC02 --- Chọn màu để gợi ý danh sách phát \\
\hline
\endhead
\textbf{Tác nhân} & Người nghe (chính); kho nhạc và đặc trưng đã tính sẵn (phụ) \\ \hline
\textbf{Mô tả} & Người nghe chọn một màu (hoặc hai màu để tạo ``hành trình cảm xúc''); hệ thống trả về danh sách bài hát có cảm xúc tương ứng và nghe đồng nhất. \\ \hline
\textbf{Tiền điều kiện} & Các tọa độ cảm xúc của kho nhạc và bảng ánh xạ màu sang Valence--Arousal đã sẵn sàng. \\ \hline
\textbf{Luồng sự kiện chính} & 1. Người nghe chọn một màu trong bảng 12 màu cơ bản. \newline 2. Hệ thống ánh xạ màu sang tọa độ cảm xúc (Valence--Arousal). \newline 3. Hệ thống chấm điểm độ gần cảm xúc trong không gian phân vị của kho nhạc, ưu tiên cụm có độ nhất quán âm thanh cao. \newline 4. Hệ thống đa dạng hóa và sắp xếp kết quả thành chuỗi chuyển cảm xúc mượt. \newline 5. Hệ thống trả về danh sách danh sách phát kèm thông tin ánh xạ màu$\rightarrow$cảm xúc. \\ \hline
\textbf{Luồng thay thế / ngoại lệ} & 1a. Người nghe chọn \textit{hai} màu: hệ thống tạo hành trình từ tâm trạng màu A sang màu B theo nguyên lý đồng cảm (iso-principle). \newline 1b. Chọn từ ba màu trở lên: hệ thống chỉ giữ hai màu đầu để tránh ``nhiễu tâm trạng''. \newline 2a. Mã màu không hợp lệ (không phải dạng \texttt{\#RRGGBB}): hệ thống từ chối yêu cầu và báo lỗi đầu vào. \\ \hline
\textbf{Hậu điều kiện} & Một danh sách phát khớp cảm xúc với màu đã chọn, có độ nhất quán âm thanh, được hiển thị cho người nghe. \\ \hline
\end{longtable}

Để minh họa rõ hai luồng (một màu và hai màu) của UC02, Hình~\ref{fig:activity-uc02} trình bày biểu đồ hoạt động tương ứng. Điểm rẽ nhánh nằm ở số màu người dùng chọn: với một màu, hệ thống đi theo nhánh truy hồi và tinh chỉnh thông thường; với hai màu, hệ thống bổ sung bước sinh hành trình cảm xúc theo nguyên lý đồng cảm trước khi trả về.

\begin{figure}[tbp]
\centering
\begin{tikzpicture}[
   term/.style={draw, rounded corners=9pt, fill=black!5, align=center, text width=3.2cm, minimum height=0.6cm, font=\small, inner sep=3pt},
   act/.style={draw, align=center, text width=3.6cm, minimum height=0.7cm, font=\small, inner sep=3pt},
   dec/.style={draw, diamond, aspect=2.2, align=center, inner sep=1pt, fill=blue!5, font=\small, text width=1.3cm},
   arr/.style={-{Stealth[length=2mm]}, thick}]

  \node[term] (s) at (0,0) {Người nghe chọn màu};
  \node[dec] (d) at (0,-1.6) {Số màu?};
  % one-colour branch (left, x=-3.8)
  \node[act] (map1) at (-3.8,-3.3) {Ánh xạ màu $\rightarrow$ V-A};
  \node[act] (cdf) at (-3.8,-4.8) {Đích phân vị (CDF)};
  \node[act] (rbf) at (-3.8,-6.3) {Chấm điểm RBF dị hướng};
  % two-colour branch (right, x=3.8)
  \node[act] (map2) at (3.8,-3.3) {Ánh xạ 2 màu $\rightarrow$ V-A};
  \node[act] (iso) at (3.8,-4.8) {Sinh hành trình (iso-principle)};
  \node[act] (sample) at (3.8,-6.3) {Lấy mẫu điểm dừng};
  % merge (center)
  \node[act, text width=5.6cm] (coh) at (0,-8.2) {Chọn cụm nhất quán âm thanh + MMR};
  \node[act, text width=5.6cm] (ret) at (0,-9.9) {Trả về danh sách phát + ánh xạ màu$\rightarrow$cảm xúc};
  \node[term] (e) at (0,-11.4) {Kết thúc};

  \draw[arr] (s) -- (d);
  \draw[arr] (d) -| (map1) node[near start, above, font=\scriptsize] {1 màu};
  \draw[arr] (d) -| (map2) node[near start, above, font=\scriptsize] {2 màu};
  \draw[arr] (map1) -- (cdf);
  \draw[arr] (cdf) -- (rbf);
  \draw[arr] (map2) -- (iso);
  \draw[arr] (iso) -- (sample);
  \draw[arr] (rbf) |- (coh);
  \draw[arr] (sample) |- (coh);
  \draw[arr] (coh) -- (ret);
  \draw[arr] (ret) -- (e);
\end{tikzpicture}
\caption{Biểu đồ hoạt động của UC02 --- gợi ý theo màu (rẽ nhánh một màu / hai màu)}
\label{fig:activity-uc02}
\end{figure}

\begin{longtable}{|p{3.0cm}|p{11.0cm}|}
\caption{Đặc tả use case UC03 --- Xem danh sách bài gợi ý}\label{table:uc03}\\
\hline
\textbf{Mã / Tên} & UC03 --- Xem danh sách bài gợi ý \\
\hline
\endfirsthead
\hline
\textbf{Mã / Tên} & UC03 --- Xem danh sách bài gợi ý \\
\hline
\endhead
\textbf{Tác nhân} & Người nghe \\ \hline
\textbf{Mô tả} & Hiển thị danh sách kết quả của UC01 hoặc UC02 để người nghe xem, phát và sắp xếp lại. \\ \hline
\textbf{Tiền điều kiện} & Một use case gợi ý (UC01/UC02) đã trả về kết quả. \\ \hline
\textbf{Luồng sự kiện chính} & 1. Hệ thống hiển thị từng bài kèm ảnh bìa (nhuộm màu theo tâm trạng), tên bài, nghệ sĩ. \newline 2. Người nghe có thể phát một bài, phát toàn bộ danh sách, hoặc kéo--thả để sắp xếp lại. \newline 3. Khi gần hết danh sách ở chế độ đài, hệ thống tự nối thêm bài (gọi lại UC01/UC02 với danh sách loại trừ). \\ \hline
\textbf{Luồng thay thế / ngoại lệ} & 1a. Danh sách rỗng: hệ thống hiển thị thông báo không có kết quả phù hợp. \newline 2a. Tệp âm thanh của bài không tồn tại: nút phát bị vô hiệu và hiển thị trạng thái không phát được. \\ \hline
\textbf{Hậu điều kiện} & Người nghe thấy được danh sách gợi ý và có thể tương tác phát nhạc; hàng đợi phát phản ánh thao tác của người nghe. \\ \hline
\end{longtable}

\begin{longtable}{|p{3.0cm}|p{11.0cm}|}
\caption{Đặc tả use case UC04 --- Xem thông tin cảm xúc bài hát}\label{table:uc04}\\
\hline
\textbf{Mã / Tên} & UC04 --- Xem thông tin cảm xúc của bài hát \\
\hline
\endfirsthead
\hline
\textbf{Mã / Tên} & UC04 --- Xem thông tin cảm xúc của bài hát \\
\hline
\endhead
\textbf{Tác nhân} & Người nghe \\ \hline
\textbf{Mô tả} & Người nghe xem thông tin cảm xúc của một bài hát: tọa độ Valence--Arousal, nhãn tâm trạng và lời bài hát. \\ \hline
\textbf{Tiền điều kiện} & Bài hát đã được gán tọa độ cảm xúc ở pha ngoại tuyến. \\ \hline
\textbf{Luồng sự kiện chính} & 1. Người nghe mở một bài hát. \newline 2. Hệ thống thể hiện tâm trạng của bài qua màu nền (suy ra từ Valence--Arousal) và một nhãn tâm trạng ngắn. \newline 3. Hệ thống cung cấp tọa độ Valence--Arousal, nhãn tâm trạng và lời bài hát qua giao diện lập trình ứng dụng (mục~\ref{section:api}). \\ \hline
\textbf{Luồng thay thế / ngoại lệ} & 3a. Bài không có lời: trường lời trả về rỗng và giao diện hiển thị ``chưa có lời''. \newline 3b. Bài thiếu một số đặc trưng âm thanh: các trường tương ứng trả về rỗng thay vì gây lỗi. \\ \hline
\textbf{Hậu điều kiện} & Thông tin cảm xúc và lời của bài hát được hiển thị/được cung cấp cho người nghe. \\ \hline
\end{longtable}

\begin{longtable}{|p{3.0cm}|p{11.0cm}|}
\caption{Đặc tả use case UC05 --- Cập nhật kho nhạc}\label{table:uc05}\\
\hline
\textbf{Mã / Tên} & UC05 --- Cập nhật kho nhạc \\
\hline
\endfirsthead
\hline
\textbf{Mã / Tên} & UC05 --- Cập nhật kho nhạc \\
\hline
\endhead
\textbf{Tác nhân} & Người vận hành / quản trị dữ liệu (chính); kho nhạc (phụ) \\ \hline
\textbf{Mô tả} & Người vận hành thu thập bài hát mới, lọc theo ngôn ngữ, tải tệp âm thanh và lời, lọc trùng/cover, rồi cập nhật vào kho nhạc. \\ \hline
\textbf{Tiền điều kiện} & Có quyền vận hành; có danh sách nghệ sĩ hạt giống và khóa truy cập các nguồn dữ liệu. \\ \hline
\textbf{Luồng sự kiện chính} & 1. Người vận hành chạy quy trình thu thập (khám phá nghệ sĩ, lấy danh sách bài). \newline 2. Hệ thống lọc bài tiếng Việt, tải tệp âm thanh và lời bài hát. \newline 3. Hệ thống phát hiện và lập danh sách bản trùng/cover. \newline 4. Kho nhạc được cập nhật với các bài hợp lệ và bảng loại trừ cover. \\ \hline
\textbf{Luồng thay thế / ngoại lệ} & 2a. Không tải được tệp âm thanh sau các mức dự phòng: bài bị bỏ qua và ghi nhật ký. \newline 2b. Không tìm được lời: bài vẫn được giữ, đánh dấu thiếu lời để pha ngoại tuyến xử lý dự phòng. \\ \hline
\textbf{Hậu điều kiện} & Kho nhạc có thêm bài mới đã lọc trùng/cover; trạng thái này là đầu vào cho UC06. \\ \hline
\end{longtable}

\begin{longtable}{|p{3.0cm}|p{11.0cm}|}
\caption{Đặc tả use case UC06 --- Chạy lại pipeline trích đặc trưng}\label{table:uc06}\\
\hline
\textbf{Mã / Tên} & UC06 --- Chạy lại pipeline trích đặc trưng \\
\hline
\endfirsthead
\hline
\textbf{Mã / Tên} & UC06 --- Chạy lại pipeline trích đặc trưng \\
\hline
\endhead
\textbf{Tác nhân} & Người vận hành (chính); hệ thống xử lý đặc trưng offline (phụ) \\ \hline
\textbf{Mô tả} & Người vận hành kích hoạt pha ngoại tuyến để trích embedding âm thanh, tính tọa độ cảm xúc và sinh các ma trận biểu diễn cho kho nhạc. \\ \hline
\textbf{Tiền điều kiện} & Kho nhạc đã cập nhật (UC05); các mô hình đóng băng và đầu dò tuyến tính đã sẵn sàng. \\ \hline
\textbf{Luồng sự kiện chính} & 1. Người vận hành chạy lần lượt các bước: trích embedding âm thanh MuQ, mã hóa lời, trích nhịp độ sạch, phát hiện cover. \newline 2. Hệ thống tính tọa độ Valence--Arousal cho từng bài và lưu ra tệp đặc trưng. \newline 3. Hệ thống đóng gói một bản phát hành phục vụ và kiểm tra đồng bộ kích thước các tệp đặc trưng. \\ \hline
\textbf{Luồng thay thế / ngoại lệ} & 1a. Quá trình trích bị gián đoạn: hệ thống hỗ trợ chạy tiếp từ điểm dừng (checkpoint). \newline 3a. Số dòng đặc trưng không khớp số bài: bước đóng gói báo lỗi và dừng để tránh phục vụ dữ liệu lệch. \\ \hline
\textbf{Hậu điều kiện} & Bộ đặc trưng mới (embedding âm thanh, tọa độ cảm xúc, embedding lời, chỉ mục cover) sẵn sàng để pha phục vụ nạp lên. \\ \hline
\end{longtable}

\section{Yêu cầu phi chức năng}
\label{section:2.4}
Ngoài các yêu cầu chức năng, hệ thống phải đáp ứng một tập yêu cầu phi chức năng (non-functional requirements) mang tính ràng buộc xuyên suốt thiết kế, tổng hợp trong Bảng~\ref{table:nfr}. Khác với nhiều ứng dụng thông thường nơi yêu cầu phi chức năng chủ yếu là hiệu năng và bảo mật, ở đồ án này nhóm yêu cầu \textit{tính đúng đắn khoa học} và \textit{khả năng kiểm chứng} được đặt lên hàng đầu, vì sản phẩm cốt lõi là một quy trình khoa học chứ không chỉ là phần mềm.

\begin{longtable}{|p{1.4cm}|p{3.2cm}|p{4.2cm}|p{4.0cm}|}
\caption{Danh sách yêu cầu phi chức năng}\label{table:nfr}\\
\hline
\textbf{Mã} & \textbf{Nhóm} & \textbf{Yêu cầu} & \textbf{Cách kiểm chứng} \\
\hline
\endfirsthead
\hline
\textbf{Mã} & \textbf{Nhóm} & \textbf{Yêu cầu} & \textbf{Cách kiểm chứng} \\
\hline
\endhead
NFR01 & Đúng đắn khoa học & Mọi tín hiệu/trọng số có nguồn từ tài liệu công bố hoặc tối ưu trên dữ liệu công khai & Truy nguồn trích dẫn; tái lập nhãn \\
\hline
NFR02 & Ràng buộc dữ liệu/mô hình & Không dùng dữ liệu người dùng; chỉ mô hình đóng băng + đầu dò tuyến tính & Rà soát mã nguồn; không có bước fine-tune \\
\hline
NFR03 & Hiệu năng & Truy vấn lúc phục vụ ở mức vài mili-giây & Đo trực tiếp (mục~\ref{subsec:latency}) \\
\hline
NFR04 & Tính tất định & Kết quả lặp lại được (chọn tham lam, tie-break ổn định đa nền tảng) & Chạy lại cho cùng kết quả \\
\hline
NFR05 & Khả năng kiểm chứng & Mọi độ đo có quy trình đánh giá tái lập được & Bộ công cụ eval chạy lại được \\
\hline
NFR06 & Bảo trì / nâng cấp & Thay một thành phần phải vượt cổng không thụt lùi trên độ đo cuối & So sánh trước/sau khi thay \\
\hline
NFR07 & Vững trước lỗi biên & Đầu vào sai/thiếu được xử lý tường minh, không gây lỗi máy chủ & Kiểm thử hợp đồng API (mục~\ref{section:functest}) \\
\hline
NFR08 & Khả năng mở rộng & Thêm bài mới không phá vỡ pipeline; phủ lời thiếu có cơ chế dự phòng & Cờ \texttt{has\_lyrics} + lùi về đặc trưng âm thanh \\
\hline
\end{longtable}

Trong các yêu cầu trên, NFR06 --- \textit{cổng không thụt lùi} (no-regression gate) --- là một yêu cầu đặc thù và quan trọng của đồ án, nên cần làm rõ. Mỗi khi một thành phần học máy được cân nhắc thay thế (ví dụ đổi mô hình embedding lời, đổi trục biểu diễn âm thanh), thành phần mới chỉ được chấp nhận nếu nó \textit{không làm xấu} độ đo cuối của hệ thống so với cấu hình hiện hành, sau khi đã được tinh chỉnh siêu tham số riêng. Ràng buộc này biến quá trình phát triển thành một chuỗi thay đổi được kiểm soát: không có cải tiến nào được đưa vào chỉ vì ``trông có vẻ tốt hơn trên giấy''. Cơ chế này được vận dụng trực tiếp trong khảo sát so sánh mô hình ở Chương~\ref{chapter:Evaluation} và được phân tích như một đóng góp phương pháp luận ở Chương~\ref{chapter:SolutionAndContribution}.

Chương này đã khảo sát hiện trạng, xác lập mô hình tác nhân, các use case chính kèm đặc tả chi tiết, cùng các yêu cầu chức năng và phi chức năng. Các nền tảng lý thuyết và công nghệ cần thiết để đáp ứng những yêu cầu này được trình bày trong chương tiếp theo.

\end{document}
